%!TEX root = ../Chimica.tex
\documentclass[11pt,oneside]{scrbook}

% ---------- LINGUA & CODIFICA ----------
\usepackage[italian]{babel}
\usepackage[T1]{fontenc}
\usepackage[utf8]{inputenc}
\usepackage{csquotes}

% ---------- TIPOGRAFIA CONSIGLIATA ----------
\usepackage{lmodern}     % font semplice e diffuso
\usepackage{microtype}   % micro-tipografia, migliore resa

% ---------- FIGURE IN TIKZ -----------

\usepackage{tikz}
\usetikzlibrary{arrows.meta, patterns, decorations.markings, decorations.pathmorphing, calc, positioning}
\usepackage{amsmath, amssymb, bm}
\usepackage{circuitikz}

% --- stili e macro (nel preambolo) ---
\tikzset{
  boccola/.style={circle,draw,thick,fill=white,inner sep=0pt,minimum size=5pt},
  cavo/.style={thick,rounded corners},
  morsetto/.style={font=\footnotesize,above=3pt},
  etichetta/.style={font=\small}
}
% quattro boccole equispaziate: b1 b2 = avv. interno, b3 b4 = avv. esterno
\newcommand{\boccole}{%
  \foreach \i in {1,2,3,4} {\node[boccola] (b\i) at (1.5*\i,0) {};}%
}
% cavo tra due boccole; argomento opzionale = looseness
\newcommand{\cavo}[3][1.3]{\draw[cavo] (#2) to[out=-70,in=-110,looseness=#1] (#3);}

% Stile coerente con le immagini a colori
\definecolor{posred}{RGB}{239,130,140}   % rosa-rosso per cariche positive
\definecolor{negblue}{RGB}{120,190,230}  % azzurro per cariche negative
\definecolor{ered}{RGB}{227,35,35}       % rosso vivo per i vettori E
\definecolor{surfblue}{RGB}{205,230,245} % azzurro chiaro per la superficie
\definecolor{gray60}{gray}{0.55}
\definecolor{blobpink}{RGB}{245,205,190} % colore del corpo esteso
\definecolor{cubecol}{RGB}{235,190,170}
\definecolor{charge}{RGB}{245,175,120}      % arancione chiaro delle cariche
\definecolor{matpink}{RGB}{245,210,195}     % superficie superiore del materiale
\definecolor{matpinkdk}{RGB}{225,175,155}   % fianco laterale
\definecolor{loopblue}{RGB}{90,130,200}     % anelli di corrente
\definecolor{mred}{RGB}{227,35,35}          % vettore M
\definecolor{surrfblue}{RGB}{190,215,235}   % regione chiusa C (riempimento)
\definecolor{surfedge}{RGB}{80,120,170}     % bordo della curva C
\definecolor{ghost}{RGB}{140,140,140}       % traslazione (tratteggiata)
\definecolor{vred}{RGB}{227,35,35}          % vettori v



% ---------- GEOMETRIA PAGINA ----------
% Scegli *una* delle due geometrie qui sotto.
% (1) Impostazioni consigliate per A4 (Italia) con ampio margine esterno:
\usepackage[
  a4paper,
  top=24mm,
  bottom=28mm,
  inner=28mm,   % margine interno (verso rilegatura)
  outer=50mm,   % margine esterno: ampio per le note a margine
  heightrounded
]{geometry}
\setlength{\headheight}{17pt} % evita sovrapposizioni con l'header
\setlength{\headsep}{6mm}     % distanza tra header e corpo


% (2) Se vuoi replicare il formato del PDF caricato (US Letter 8.5x11"):
% \usepackage[
%   letterpaper,
%   top=25mm,
%   bottom=30mm,
%   inner=30mm,
%   outer=55mm,
%   heightrounded
% ]{geometry}

% ---------- HEADER/FOOTER ----------
\usepackage{fancyhdr}
\pagestyle{fancy}
\fancyhf{} % pulisci

% intestazione: solo titolo del capitolo (sinistra), pagina a destra
\fancyhead[L]{\nouppercase{\leftmark}}
\fancyfoot[R]{\thepage}

\renewcommand{\headrulewidth}{0.4pt}
\renewcommand{\footrulewidth}{0pt}

\fancypagestyle{plain}{%
  \fancyhf{}
  \fancyfoot[R]{\thepage}
  \renewcommand{\headrulewidth}{0pt}
}

% ---------- STRUMENTI DI STRUTTURA ----------
\setcounter{secnumdepth}{3}
\setcounter{tocdepth}{2}
\usepackage{titlesec}

% --- Titoli puliti ---
% Capitolo
\titleformat{\chapter}[display]
  {\normalfont\huge\bfseries}
  {\thechapter}{1ex}{}
\titlespacing*{\chapter}{0pt}{-0.5em}{1.5ex}

% Capitolo
\titleformat{\chapter}[display]
  {\normalfont\huge\bfseries}
  {\thechapter}{1ex}{}
\titlespacing*{\chapter}{0pt}{1.5ex}{1.5ex}

% Sezione
\titleformat{\section}
  {\normalfont\Large\bfseries}
  {\thesection}{1em}{}

% Sottosezione
\titleformat{\subsection}
  {\normalfont\large\bfseries}
  {\thesubsection}{0.8em}{}

% SOTTO-sottosezione  <--- AGGIUNGI QUESTO BLOCCO
\titleformat{\subsubsection}
  {\normalfont\normalsize\bfseries}
  {\thesubsubsection}{0.8em}{}
\titlespacing*{\subsubsection}{0pt}{0.8ex}{0.4ex}

% Paragrafo (run-in: titolo in grassetto sulla stessa riga del testo)
\titleformat{\paragraph}[runin]
  {\normalfont\normalsize\bfseries}
  {}{0pt}{}[.]
\titlespacing*{\paragraph}{0pt}{1ex}{0.5em}

% ---------- PAGINA SEPARATRICE ----------
\newcommand{\partpage}[1]{%
  \cleardoublepage          % vai a pagina nuova (pari se oneside)
  \thispagestyle{empty}     % niente header/footer
  \null\vfill               % spazio verticale elastico sopra
  \begin{center}
    {\Huge\bfseries #1}
  \end{center}
  \vfill\null               % spazio verticale elastico sotto
  \cleardoublepage          % lascia la pagina successiva libera
}

% ---------- AMBIENTI UTILI ----------
\usepackage{amsmath,amssymb,amsthm}
\usepackage{graphicx}
\usepackage[hidelinks]{hyperref}
\usepackage{xcolor}
\usepackage{tcolorbox}
\usepackage{subcaption}
\usepackage{booktabs}
\usepackage{multirow}
\usepackage[version=4]{mhchem}   % per \ce{}
\usepackage{siunitx}             % per \SI{} e \si{}


% ---------- NOTE A MARGINE (parole-chiave/sidenote) ----------
% Definiamo un comando comodo e graficamente coerente:
\newcommand{\mkeyword}[1]{%
  \marginpar{\raggedright\footnotesize\emph{#1}}%
}

% Se preferisci le note a sinistra, attiva la riga seguente:
% \reversemarginpar

% ---------- PARAGRAFI E SPAZI ----------
\setlength{\parindent}{0.8em}  % rientro tradizionale
\setlength{\parskip}{0.25em}   % piccolo spazio extra tra paragrafi

%----------- COMANDI AGGIUNTIVI -----------
% --- Comando per vettori colonna di lunghezza variabile ---

\newcommand{\colvec}[1]{%
  \begin{pmatrix}
    #1
  \end{pmatrix}%
}

\renewcommand{\Re}{\operatorname{Re}}
\renewcommand{\Im}{\operatorname{Im}}

\newcommand{\FIG}[1]{
  \begin{figure}[h]
    \centering
    \includegraphics[width=0.4\linewidth]{figure/image#1.png}
    \caption{}
    \label{}
  \end{figure}
}

\newcommand{\eps}{\varepsilon}

\theoremstyle{definition}
\newtheorem{definizione}{Definizione}[chapter]
\newtheorem{esempio}[definizione]{Esempio}

\theoremstyle{remark}
\newtheorem{osservazione}[definizione]{Osservazione}

\theoremstyle{plain}
\newtheorem{teorema}[definizione]{Teorema}
\newtheorem{proposizione}[definizione]{Proposizione}

% % ---------- BIBLIOGRAFIA ----------
% \usepackage[backend=biber,style=numeric,sorting=nyt]{biblatex}
% \addbibresource{chapters/bibliografia.bib}

% ---------- FRONTESPIZIO ----------
\title{\textbf{Esperienze di Laboratorio 2}}
\author{Roberto Riccucci}
\date{}

\publishers{
  Email: \text{roberto.riccucci1304@gmail.com} \\
  Versione: v0.1
}

